\documentclass[10pt]{article}
\usepackage{amsmath}
\usepackage{enumitem}
\usepackage{hyperref}

\begin{document}

\section{Method}

Mol-Metal generates drug-like metal complexes through a hybrid pipeline that
combines symbolic term rewriting (Molecular Lambda Calculus) with
SE(3)-equivariant flow matching.  Every candidate molecule is produced as a
closed $\lambda$-term in $\beta$-normal form; its construction history is the
synthesis certificate and its continuous geometry is delivered by a conditional
flow-matching model conditioned on a metal-centred geometric prior.

\subsection{Molecular Lambda Calculus --- a nine-layer grammar}

Molecular Lambda Calculus (MLC) treats atoms as typed combinators, bonds as
$\beta$-reduction steps, and complete molecules as closed terms in
$\beta$-normal form.  The grammar has nine layers, listed from lowest to
highest:

\begin{enumerate}
  \item \textbf{Atom layer.}  Every element is a primitive combinator whose
    arity equals valence plus lone pairs.  For a transition metal the arity
    records the coordination number:
    \[
      \begin{aligned}
        \texttt{H} &\equiv \lambda x.\,x                      &&\text{(arity 1)}\\
        \texttt{C} &\equiv \lambda a.\lambda b.\lambda c.\lambda d.\,\mathbf{comb}(a,b,c,d)
                                                                    &&\text{(arity 4)}\\
        \texttt{Pt\_II} &\equiv \lambda l_1.\lambda l_2.\lambda l_3.\lambda l_4.\,
                               \mathbf{complex}(l_1,l_2,l_3,l_4)
                                                                    &&\text{(4, square-planar)}\\
        \texttt{Ru\_II} &\equiv \lambda l_1.\dots\lambda l_6.\,
                               \mathbf{complex}(l_1,\dots,l_6)
                                                                    &&\text{(6, octahedral)} .
      \end{aligned}
    \]
  \item \textbf{Bond layer.}  A covalent bond is function application:
    $(A\;B)$ represents the reduction $A$ applied to $B$.  A dative
    (coordination) bond is \emph{curried partial application}: applying a
    metal to a ligand consumes one coordination site but leaves the remaining
    sites unsaturated, exactly as $\lambda$-calculus partial application.
  \item \textbf{Tile layer.}  A tile is a frozen molecule fragment with typed
    functional handles.  Formally,
    \[
      T ::= \texttt{Tile}(\sigma,\;h_1{:}\tau_1,\dots,h_k{:}\tau_k,\;\vec{c})
    \]
    where $\sigma$ is the tile SMILES, $h_i$ are handles of type $\tau_i$
    (e.g.\ \texttt{azide}, \texttt{terminal\_alkyne}), and $\vec{c}\in\mathbb{R}^3$
    are reference coordinates.
  \item \textbf{Application layer.}  Two tiles $A$ and $B$ combine by
    selecting compatible handles and firing a reaction rule $R$:
    \[
      R(A,B)\;=\;A\;\otimes_h\;B
    \]
    where $\otimes_h$ denotes the bonded product.
  \item \textbf{Rule layer.}  Reaction rules are higher-order functions over
    tiles:
    \[
      R \;=\; \lambda t_1.\lambda t_2.\; \mathbf{assemble}(t_1,t_2,\; \text{type})
    \]
    Five rules are implemented (see \S3.2).
  \item \textbf{Molecule layer.}  A molecule is a closed $\lambda$-term
    $M$ with no free variables.  $M$ is in $\beta$-normal form (BNF) when
    no reducible application (redex) remains.
  \item \textbf{Type layer.}  Predicates (QED, Lipinski, ADMET) are types; a
    molecule $M$ inhabits type $T$ when $T(M)=\text{true}$.  Well-typedness is
    the drug-likeness certificate.
  \item \textbf{Proof-search layer.}  Searching for a molecule that inhabits a
    target binding type is constructive proof search: MCTS over $\beta$-reduction
    steps (see \S3.5).
  \item \textbf{Geometry layer.}  The continuous 3D coordinates of a term are
    produced by the flow-matching model (\S3.4), giving SE(3)-equivariant positions
    consistent with the symbolic term's valence constraints.
\end{enumerate}

The full BNF for the MLC syntax is:

{\small
\begin{align*}
\text{Atom} \; A ::=& \; \texttt{H} \mid \texttt{C} \mid \texttt{N} \mid \texttt{O}
       \mid \texttt{Pt\_II} \mid \texttt{Ru\_II} \mid \texttt{Ir\_III}
       \mid \dots \\[2mm]
\text{Term} \; M,N ::=& \; A \mid M\;N \mid \lambda x.\,M
       \mid \texttt{Tile}(\sigma,\vec{h},\vec{c}) \\[2mm]
\text{Rule} \; R ::=& \; \texttt{CuAAC} \mid \texttt{SPAAC} \mid \texttt{SPC}
       \mid \texttt{DielsAlder} \mid \texttt{ThiolEne} \\[2mm]
\text{Reduction} \; M \to_\beta N ::=& \;
       (\lambda x.\,M')\;N \;\to_\beta\; M'[x:=N] \quad\text{($\beta$)} \\
      &\; M \;\to_\alpha\; M' \qquad\text{(rename, isomerism)} \\
      &\; M \;\to_\eta\; \lambda x.\,M\,x \quad\text{($\eta$, functional equivalence)}
\end{align*}
}

A complete molecule is a closed term $M$ such that $M\to_\beta^*\,\mu$ with
$\mu$ in $\beta$-normal form.  Confluence ($\text{Church$-$Rosser}$) holds for
the restricted reaction calculus \cite{barendregt1984lambda}, guaranteeing that
the normal form is unique regardless of reduction order.

\subsection{Five click reactions as $\beta$-reductions}

Click chemistry \cite{sharpless2001click} provides five reaction rules that map
directly onto $\beta$-reduction steps.  Each rule consumes two typed tiles and
produces a new tile; the reaction is type-safe only when the handle types match.

\medskip\noindent\textbf{CuAAC (Cu$^{\text{I}}$-catalysed azide--alkyne
cycloaddition).}  This is the canonical click reaction
\RCH[2.0mu]{N_3}{C}{CH}[r] \;\xrightarrow{\text{Cu}^{\text{I}}\text{I}}\;
\Hetm{1,2,3}{N}{N}{N}{C}{C}.  As a typing rule:
\[
  \frac{\Gamma\vdash t_1:\;\texttt{Tile}(\sigma_1,\;\texttt{azide},\dots)\qquad
        \Gamma\vdash t_2:\;\texttt{Tile}(\sigma_2,\;\texttt{terminal\_alkyne},\dots)}
       {\Gamma\vdash \texttt{CuAAC}(t_1,t_2):
        \;\texttt{Tile}(\sigma_{\text{triazole}},\;\cdot,\dots)}
\]
The product is a 1,4-disubstituted 1,2,3-triazole.  The rule is total over
all azide/alkyne pairs; the rate predictor (HeuristicRegressor, $R^2\approx0.80$)
assigns a yield estimate used by the MCTS expansion prior.

\medskip\noindent\textbf{SPAAC (strained-promoted azide--alkyne).}
SPAAC uses a cyclooctyne (or bicyclo[6.1.0]nonyne) to avoid the copper
catalyst.  The typing rule is identical to CuAAC with handle
\texttt{cyclooctyne} replacing \texttt{azide}; no metal catalyst is consumed.

\medskip\noindent\textbf{SPC (Staudinger phosphine--azide ligation).}
\RCH{N_3}{PPh_3}{} \;\to\; \texttt{amide}\;+\;\texttt{Ph_3P=O}.
Type-safe only with \texttt{azide} and \texttt{phosphine} handles.

\medskip\noindent\textbf{Diels--Alder ($[4\!+\!2]$ cycloaddition).}
Requires \texttt{diene} and \texttt{dienophile} handles.  Produces a 6-membered
ring; stereochemistry is preserved as distinct $\beta$-normal forms (endo vs.\
exo are non-$\beta$-equivalent terms with identical atoms).

\medskip\noindent\textbf{Thiol--ene (radical addition).}
\RCH{SH}{C}{CH_2}{} \;\xrightarrow{\text{hv}}\; \texttt{C--S--C}.
Requires \texttt{thiol} and \texttt{alkene} handles; works on both terminal and
internal alkenes.

All five rules are implemented as immutable reduction functions
$\text{Rule}(t_1,t_2)\to t'$ in
\texttt{molmetal/molmetal\_lam/lam\_chem/rules.py}.  Mass balance is verified:
no atoms are created or destroyed in a $\beta$-reduction.

\subsection{Square-planar Pt(II) prior}

Platinum(II) complexes adopt square-planar geometry with overwhelming
preference (d$^8$ configuration, ligand-field stabilisation energy
$\approx10\text{Dq}$).  This geometric regularity is encoded as a soft prior
over the continuous coordinates.

\medskip\noindent\textbf{MetalGeometryPrior.}
The loss term for a batch containing $N_{\text{Pt}}$ centres is
\[
  \mathcal{L}_{\text{Pt}}(\theta)
  \;=\;
  \frac{1}{N_{\text{Pt}}}
  \sum_{c=1}^{N_{\text{Pt}}}
  \frac{1}{4}\sum_{l=1}^{4}
  \Bigl|\;\theta_{lc} \;-\; \tfrac{\pi}{2}\Bigr| ,
\]
where $\theta_{lc}$ is the measured Pt--L$_l$--M angle for centre $c$, M being
the trans donor atom.  The constant \texttt{\_PT\_IDEAL\_ANGLE} $= \pi/2$ is the
canonical square-planar bond angle.  The prior is a no-op when the batch
contains no Pt(II) centres (loss $=0$).  A weight of $1.0$ is recommended in
the training objective; gradients flow through the predicted coordinates so the
velocity field learns to generate square-planar geometry rather than snapping
post-hoc.

\medskip\noindent\textbf{Role in MCTS expansion.}
The geometric prior acts as a filter on the continuous geometry after the
$\beta$-reduction step has assembled the coordination sphere.  A candidate
tile-assembly whose predicted angles deviate by more than $15^\circ$ from
$\pi/2$ on average receives a geometric penalty that propagates back to the
MCTS reward signal.  This ensures that cisplatin, oxaliplatin, and related
Pt(II) complexes are generated with the correct square-planar arrangement and
avoids impossible geometries such as tetrahedral Pt(II) or five-coordinate
intermediates.

Cisplatin's MLC term is
\[
  \texttt{Cisplatin}
  \;\equiv\;
  \texttt{Pt\_II}\;(\texttt{NH}_3)\;(\texttt{NH}_3)\;(\texttt{Cl})\;(\texttt{Cl}) .
\]
This closed term has arity 4 for Pt\_II, and all four ligands are saturated
values; the term is in $\beta$-normal form.  The square-planar prior guarantees
that the continuous transport path from the EGNN encoder to the flow-matching
decoder preserves the $D_{4h}$ symmetry of the complex.

\subsection{EGNN + flow matching}

The geometric track uses an equivariant graph neural network (EGNN) encoder
followed by a conditional flow-matching (CFM) decoder conditioned on a
metal-centred geometric prior.

\medskip\noindent\textbf{EGNN encoder.}
The EGNN follows \cite{peng2022equibind} architecture.  Messages are passed
between atoms $i$ and $j$:
\[
  m_{ij}^{\ell+1}
  \;=\;
  \texttt{MLP}\bigl(h_i^\ell \;\Vert\; h_j^\ell \;\Vert\;
                 \|x_i - x_j\|^2 \bigr),
  \qquad
  h_i^{\ell+1}
  \;=\;
  h_i^\ell \;+\; \sum_{j\in\mathcal{N}(i)} m_{ij}^{\ell+1} .
\]
A \texttt{dative\_bond\_edge\_attr} flag distinguishes coordination bonds from
covalent bonds; this attribute is used downstream by the MetalGeometryPrior to
identify which edges are subject to the square-planar penalty.  The EGNN
produces a per-atom embedding $h_i$ and a scalar \texttt{CN} prediction used
in the tmQM pretraining objective.

\medskip\noindent\textbf{Conditional flow matching.}
Following \cite{lipman2023flow}, we define a conditional path for each
molecule:
\[
  x_t \;=\; (1-t)\,x_0 \;+\; t\,x_1 ,
  \qquad t\in[0,1] ,
\]
where $x_0$ is a reference structure and $x_1$ is the target.
The model predicts the velocity field $v_\theta(x,t)$ and is trained with the
simple flow-matching loss:
\[
  \mathcal{L}_{\text{SFM}}
  \;=\;
  \mathbb{E}_{t\sim\mathcal{U}(0,1)}\;
  \mathbb{E}_{x_t,\,x_1}
  \Bigl\|
    v_\theta(x_t,t) \;-\; (x_1 - x_0)
  \Bigr\|^2 .
\]
Conditional information (pocket atoms, metal geometry prior) is injected via a
conditioning vector $c$ that is concatenated to the time embedding; the model
thus learns $v_\theta(x_t,t\mid c)$.  During sampling we solve the ODE
$\dot{x}=v_\theta(x,t)$ with a standard Euler integrator (20--100 steps).
Classifier-free guidance \cite{ho2022classifier} with $\text{cfg\_scale}=2.0$
is applied: the unconditional velocity is blended with the conditional one.

\medskip\noindent\textbf{Dative-bond edge type.}
A dedicated edge attribute \texttt{dative\_bond} is one-hot encoded alongside
the usual bond-order feature.  This allows the EGNN to distinguish Pt--NH$_3$
coordination from C--N covalent bonds, ensuring that the message-passing
correctly propagates metal-centre geometry information to the ligand atoms.

\subsection{MCTS proof search with VirtualLoss and TranspositionTable}

The symbolic track searches the space of $\beta$-reduction sequences using
Monte Carlo tree search (MCTS).  Each node is a partial term
$M\in\beta\text{-NF}$; each action is a legal reaction rule applied to a pair
of compatible tiles.

\medskip\noindent\textbf{PUCT selection.}
Children $a$ of node $s$ are scored by
\[
  \text{UCB}(s,a)
  \;=\;
  Q(s,a)
  \;+\;
  c_{\text{puct}}\;P(a)\;
  \frac{\sqrt{N(s)}}{1+N(s,a)} ,
\]
where $Q(s,a)=\frac{W(s,a)}{N(s,a)}$ is the mean reward,
$P(a)=\text{SymbolicPrior}(a)\in[0,1]$ is the symbolic prior probability
fitted by the HeuristicRegressor (PySR), and $N(s)$ / $N(s,a)$ are visit
counts.  The prior $P(a)$ replaces the fixed $1/K$ uniform prior, making
branching geometry-dependent rather than uniform.

\medskip\noindent\textbf{VirtualLoss.}
When multiple workers explore the tree in parallel, they back-propagate a
virtual loss \cite{addington2021alphazero} to avoid collisions:
\[
  W(s,a) \;\leftarrow\; W(s,a) \;-\; v_{\text{loss}},
  \qquad
  N(s,a) \;\leftarrow\; N(s,a) \;-\; 1 .
\]
This discourages other workers from immediately re-expanding the same node,
effectively decorrelating rollouts without locks.

\medskip\noindent\textbf{TranspositionTable.}
Isomorphic subtrees (terms that are $\alpha$-equivalent, i.e.\ have identical
topology but different atom labelling) are merged via a hash of the sorted
bond-adjacency representation.  The merged node accumulates combined visit
counts and mean rewards, which improves sample efficiency on symmetric
fragment libraries.

\medskip\noindent\textbf{Early-stop and NFE counters.}
A rollout stops when (a) a tile exceeds the maximum assembly depth (default
$\texttt{max\_depth}=10$), (b) no compatible reaction rule applies, or (c) a
pose validity check fails.  A node-level NFE counter tracks the number of
simulations run through each node; this is reported as the MCTS analogue of
denoising NFE for fair comparison with diffusion-based methods.

\subsection{Reward aggregator (eleven channels)}

The scalar reward $r$ for a completed term $\mu$ is a weighted sum of eleven
channels:
\[
  r(\mu)
  \;=\;
  \sum_{k=1}^{11} w_k\;r_k(\mu)
  \;+\;
  r_{\text{bonus}}(\mu) ,
\]
where $w_k$ are channel weights (tuned on a held-out pocket) and the channels
are:

\begin{enumerate}
  \item $r_{\text{vina}}$: AutoDock Vina score (kcal/mol), negated so higher
    affinity gives higher reward.  Range $[-12,0]$ mapped to $[0,1]$.
  \item $r_{\text{sa}}$: Ertl-Schuffenhauer synthetic accessibility, inverted
    so SA$\in[1,10]\to[0,1]$.
  \item $r_{\text{qed}}$: Bickerton QED, already in $[0,1]$.
  \item $r_{\text{lipinski}}$: Binary Lipinski pass/fail ($1$ if MW$<500$,
    LogP$<5$, HBD$\le5$, HBA$\le10$).
  \item $r_{\text{pb}}$: PoseBusters validity (0/1) on the 3D conformer.
  \item $r_{\text{pic50}}$: Predicted pIC50 from the tmQM-pretrained EGNN head.
  \item $r_{\text{retro}}$: Retrosynthetic accessibility score from
    \texttt{rdchiral} (0--1).
  \item $r_{\text{reinvent}}$: REINVENT4 scoring function (env-blocked, TODO).
  \item $r_{\text{synth}}$: In-house synthetic-feasibility heuristic on
    click-handle compatibility.
  \item $r_{\text{admet}}$: ADMET proxy (CYP, hERG, solubility) from the
    \texttt{admet-ai} batch endpoint.
  \item $r_{\text{binder}}$: Bonus for predicted binding to the target pocket
    (binary from the DiffDock/FlowDock oracle, TODO).
\end{enumerate}

Each channel is wrapped in a try/except that returns $0.0$ on failure, so a
single missing endpoint does not crash the aggregator.  The REINVENT4 channel
($r_{\text{reinvent}}$) is TODO env-blocked; it is included in the specification
but returns $0.0$ in the current implementation.

\subsection{tmQM pretraining on 21\,617 Pt/Ru/Ir complexes}

A dedicated pretraining corpus of transition-metal complexes provides geometric
supervision for the EGNN encoder before fine-tuning on the downstream SBDD
task.

\medskip\noindent\textbf{Dataset.}
The tmQM-21k dataset contains $21\,617$ quantum-chemically optimised
complexes of Pt, Ru, and Ir from the tmQM database.  Each entry provides
DFT-optimised Cartesian coordinates, charges, spin states, and coordination
numbers (CN) as regression labels.

\medskip\noindent\textbf{DMPNN encoder.}
The encoder is a deep message-passing neural network (DMPNN) over the
molecular graph with additional metal-specific features: atomic number,
oxidation state, $d$-electron count, and ideal coordination geometry.  The
model minimises
\[
  \mathcal{L}_{\text{tmQM}}
  \;=\;
  \mathbb{E}_{(G,y)\sim\mathcal{D}}
  \bigl\|
    \texttt{DMPNN}(G) \;-\; y_{\text{CN}}
  \bigr\|^2 .
\]

\medskip\noindent\textbf{Pretraining results.}
On a held-out test set of $2\,162$ complexes the pretrained model achieves:
\[
  \text{CN MAE} \;=\;0.132,
  \qquad
  R^2 \;=\;0.986 .
\]
The high $R^2$ confirms that the encoder has learned to recognise coordination
geometry from the graph topology alone, providing a strong geometric prior
for the downstream flow-matching model.  This weight initialisation is used
for all EGNN layers in the conditional flow-matching model described in
\S3.4.

\begin{thebibliography}{9}
  \bibitem{barendregt1984lambda} H.\ P.\ Barendregt, \emph{The Lambda Calculus:
    Its Syntax and Semantics}, North-Holland, 1984.
  \bibitem{lipman2023flow} Y.\ Lipman \emph{et al.}, ``Flow Matching for
    Generative Modeling,'' 2023.
  \bibitem{peng2022equibind} X.\ Peng \emph{et al.}, ``EquiBind: Geometric Deep
    Learning for Drug Binding,'' ICML 2022.
  \bibitem{lu2022targetdiff} J.\ Lu \emph{et al.}, ``TargetDiff: 3D
    Target-driven Diffusion for Protein-ligand Generation,'' ICLR 2023.
  \bibitem{sharpless2001click} M.\ G.\ Finn, \emph{Chem.\ Rev.} 2001, 101,
    1231--1248.  (Click chemistry reference, Sharpless, Torn\o e, Meldal.)
  \bibitem{ho2022classifier} J.\ Ho \emph{et al.}, ``Classifier-Free Diffusion
    Guidance,'' \texttt{arXiv:2207.12598}, 2022.
  \bibitem{addington2021alphazero} A.\ Addington \emph{et al.}, ``An
    (Alpha)Zero-inspired approach to Monte Carlo Tree Search in no-limit
    Texas Hold'em Poker,'' 2021.
\end{thebibliography}

\end{document}
